\section*{Capítulo \ref{ch:g7:chemical-reactions} --- Las reacciones químicas: reactivos y productos}

\begin{solution}{exo:g7:chemical-reactions:1}
Físicas: el hielo que se funde, el azúcar que se \omterm{def:g2:mixing-and-dissolving:dissolve}{disuelve}. Químicas:
una cerilla que \omterm{def:g4:what-a-fire-needs:burning}{arde}, el pan que se tuesta, un clavo que se oxida.
\end{solution}

\begin{solution}{exo:g7:chemical-reactions:2}
Los \omterm{def:g7:chemical-reactions:reactant}{reactivos} son las especies que se consumen; los \omterm{def:g7:chemical-reactions:reactant}{productos} son las
especies nuevas que se forman.
\end{solution}

\begin{solution}{exo:g7:chemical-reactions:3}
carbono $+$ oxígeno $\longrightarrow$ dióxido de carbono.
\end{solution}

\begin{solution}{exo:g7:chemical-reactions:4}
\omterm{def:g7:chemical-reactions:reactant}{Reactivos}: el hierro y el oxígeno. \omterm{def:g7:chemical-reactions:reactant}{Producto}: el óxido de hierro.
\end{solution}

\begin{solution}{exo:g7:chemical-reactions:5}
El vaho muestra agua; el agua de cal turbia muestra dióxido de carbono.
\end{solution}

\begin{solution}{exo:g7:chemical-reactions:6}
cinc $+$ ácido clorhídrico $\longrightarrow$ cloruro de cinc $+$ hidrógeno.
\end{solution}

\begin{solution}{exo:g7:chemical-reactions:7}
Antes: 4 (en la \omterm{def:g7:atoms-and-molecules:molecule}{molécula} de metano). Después: $2 + 2 = 4$ (dos en cada
\omterm{def:g7:atoms-and-molecules:molecule}{molécula} de agua).
\end{solution}

\begin{solution}{exo:g7:chemical-reactions:8}
No aparece ninguna especie nueva: la sal sigue ahí, repartida por el
agua, y se recupera \omterm{def:g3:separating-mixtures:evaporate}{evaporando} el agua. Es una \omterm{def:g7:chemical-reactions:physical-transformation}{transformación física}.
\end{solution}

\begin{solution}{exo:g7:chemical-reactions:9}
La mayor parte de la madera ha reaccionado con el oxígeno y ha dado
\omterm{def:g7:chemical-reactions:reactant}{productos} que son gases, dióxido de carbono y vapor de agua, que se van
al \omterm{def:g4:air-a-mixture-of-gases:air}{aire}. Solo queda la ceniza.
\end{solution}

\begin{solution}{exo:g7:chemical-reactions:10}
hidrógeno $+$ oxígeno $\longrightarrow$ agua. El \omterm{def:g7:chemical-reactions:reactant}{producto} vuelve azul el
sulfato de cobre anhidro.
\end{solution}

\begin{solution}{exo:g7:chemical-reactions:11}
Antes: $2 \times 2 = 4$ \omterm{def:g7:atoms-and-molecules:atom}{átomos} de hidrógeno y 2 \omterm{def:g7:atoms-and-molecules:atom}{átomos} de oxígeno.
Después: dos \omterm{def:g7:atoms-and-molecules:molecule}{moléculas} de agua contienen $2 \times 2 = 4$ \omterm{def:g7:atoms-and-molecules:atom}{átomos} de
hidrógeno y $2 \times 1 =
2$ \omterm{def:g7:atoms-and-molecules:atom}{átomos} de oxígeno. Los mismos \omterm{def:g7:atoms-and-molecules:atom}{átomos}, en el mismo número.
\end{solution}

\begin{solution}{exo:g7:chemical-reactions:12}
Un \omterm{def:g7:atoms-and-molecules:atom}{átomo} de oxígeno no puede salir de la nada: el dibujo tiene 2 \omterm{def:g7:atoms-and-molecules:atom}{átomos}
de oxígeno antes y 3 después. Correcto: un \omterm{def:g7:atoms-and-molecules:atom}{átomo} de carbono $+$ una
\omterm{def:g7:atoms-and-molecules:molecule}{molécula} de oxígeno
$\longrightarrow$ una \omterm{def:g7:atoms-and-molecules:molecule}{molécula} de dióxido de carbono, sin que sobre
nada.
\end{solution}

\begin{solution}{pb:g7:chemical-reactions:1}
\textbf{1.} El metano y el oxígeno.

\textbf{2.} Agua.

\textbf{3.} Dióxido de carbono.

\textbf{4.} No: no se consume, y no se escribe en la \omterm{def:g7:chemical-reactions:word-equation}{ecuación con
palabras}.

\textbf{5.} metano $+$ oxígeno $\longrightarrow$ dióxido de carbono $+$
agua.

\textbf{6.} Química: desaparecen dos especies (el metano y el oxígeno) y
aparecen dos especies nuevas, detectadas mediante ensayos; además, se
desprenden calor y luz.

\textbf{7.} \ce{CH4}, \ce{O2}, \ce{CO2}, \ce{H2O}.

\textbf{8.} 1 de carbono, 4 de hidrógeno, $2 \times 2 = 4$ de oxígeno.

\textbf{9.} 1 de carbono; $2 \times 2 = 4$ de hidrógeno; $2 + 2 \times 1 = 4$
de oxígeno. Los mismos \omterm{def:g7:atoms-and-molecules:atom}{átomos} en el mismo número: solo se han
reorganizado.

\textbf{10.} $10 \times 2 = 20$ \omterm{def:g7:atoms-and-molecules:molecule}{moléculas} de oxígeno.

\textbf{11.} 10 \omterm{def:g7:atoms-and-molecules:molecule}{moléculas} de dióxido de carbono.

\textbf{12.} $10 \times 2 = 20$ \omterm{def:g7:atoms-and-molecules:molecule}{moléculas} de agua.
\end{solution}
